%%%%%%%%%%%%%%%%%%%%%%%%%%%%%%%%%%%%%%%%%%%%%%%%%%%%%%%%%%%%%%%%%%%%%%%%%%%%%%%%%%%%%%%%%
%%%%%%%%%%%%%%%%%%%%%%%%%%%%%%%%%%%%%%%% ATENÇÃO %%%%%%%%%%%%%%%%%%%%%%%%%%%%%%%%%%%%%%%%

%% Adaptado para gradução por Valentin Nakayama, março de 2024. Modificações listadas em README.txt.

%% Este template foi modificado por Aline Novais e revisado por Cláudia Maria Fortes da Silva, Maria Alice Ciocca de Oliveira e Regina de Moura para uso em Dissertações, Qualificações e Teses de alunos do Observatório do Valongo (UFRJ). Template original de Steven Gunn e Sunil Patel (mais informações abaixo).
%%% eu, Yanna, revisei conforme as recomendaçoes da biblioteca e fiz as modificações necessárias pra deixar isso no grauuu (28mar23)

%% Se você tem dúvidas de como utilizar este template em LaTeX, veja o arquivo README.txt.

%% Comentários em inglês demarcados por % são do template original. Comentários em português demarcados por %% foram feitos por Aline.

%% Caso deseje escrever o seu trabalho em inglês, comente a linha 59 (pacote "babel") e modifique as demais linhas necessárias, de acordo com os comentários.

%% Não se esqueça de modificar também o arquivo Thesis.cls com as informações do seu trabalho.

%% (Aline, 03/08/2022)

%%%%%%%%%%%%%%%%%%%%%%%%%%%%%%%%%%%%%%%%%%%%%%%%%%%%%%%%%%%%%%%%%%%%%%%%%%%%%%%%%%%%%%%%%
%%%%%%%%%%%%%%%%%%%%%%%%%%%%%%%%%%%%%%%%%%%%%%%%%%%%%%%%%%%%%%%%%%%%%%%%%%%%%%%%%%%%%%%%%


% Masters/Doctoral Thesis 
% LaTeX Template
% Version 1.43 (17/5/14)
%
% This template has been downloaded from:
% http://www.LaTeXTemplates.com
%
% Original authors:
% Steven Gunn 
% http://users.ecs.soton.ac.uk/srg/softwaretools/document/templates/
% and
% Sunil Patel
% http://www.sunilpatel.co.uk/thesis-template/
%
% License:
% CC BY-NC-SA 3.0 (http://creativecommons.org/licenses/by-nc-sa/3.0/)
%
% Note:
% Make sure to edit document variables in the Thesis.cls file


%----------------------------------------------------------------------------------------
%   PACOTES E CONFIGURAÇÕES DO DOCUMENTO
%----------------------------------------------------------------------------------------


\documentclass[11pt, oneside]{Thesis} % The default font size and one-sided printing (no margin offsets)

\graphicspath{{Imagens/}} % Specifies the directory where pictures are stored

\usepackage[round, colon]{natbib} % Use the natbib reference package - read up on this to edit the reference style; if you want text (e.g. Smith et al., 2012) for the in-text references (instead of numbers), remove 'numbers' 
\hypersetup{urlcolor=black, colorlinks=true} % Colors hyperlinks in blue - change to black if annoying
\title{\ttitle} % Defines the thesis title - don't touch this


\usepackage{longtable}
\usepackage{pdflscape}
\usepackage{array}
\usepackage{tabularx}
\usepackage{moresize}
\usepackage[brazilian]{babel}   %% comente esta linha caso seu trabalho seja em inglês
\usepackage{helvet}
\usepackage{units}
\usepackage{wrapfig}
\usepackage{subcaption}
\usepackage{indentfirst}
\usepackage{xstring}
\usepackage{amsfonts}
\usepackage{enumerate}
\usepackage[font=scriptsize,labelsep=period]{caption}
\usepackage{url}
\usepackage{booktabs}
\usepackage{afterpage}
\everymath{\displaystyle}
\newcommand{\PR}[1]{\ensuremath{\left[#1\right]}}
\newcommand{\PC}[1]{\ensuremath{\left(#1\right)}}
\newcommand{\chav}[1]{\ensuremath{\left\{#1\right\}}}
\usepackage{mdframed}
\usepackage{wasysym}
\usepackage{amsmath}
\usepackage{amssymb}
\usepackage{float}
\usepackage[absolute]{textpos}
\usepackage{hyphenat}
\usepackage{ragged2e}
\usepackage{lipsum}
\usepackage[table,x11names]{xcolor}

\AtBeginDocument{%
%\newcommand{\todayen}{\monthname \space \number \year}
\renewcommand{\today}{\ifcase \month \or Janeiro\or Fevereiro\or Março\or Abril\or Maio\or Junho\or Julho\or Agosto\or Setembro\or Outubro\or Novembro\or Dezembro\fi \space de \number \year}
}

\setcitestyle{round, aysep={,}, yysep={,}, notesep={}}

%\numberwithin{equation}{section}
%\usepackage[hmarginratio=1:1]{geometry}

\def\mean#1{\left< #1 \right>}

%\usepackage[utf8]{inputenc}
%\usepackage[T1]{fontenc}
 
\newcommand{\specialcell}[2][c]{%
  \begin{tabular}[#1]{@{}l@{}}#2\end{tabular}}


\begin{document}



%\Pre-Texto % Use roman page numbering style (i, ii, iii, iv...) for the pre-content pages

\setstretch{1.3} % Line spacing of 1.3

% Define the page headers using the FancyHdr package and set up for one-sided printing
\fancyhead{} % Clears all page headers and footers
\rhead{\thepage} % Sets the right side header to show the page number
\lhead{} % Clears the left side page header

\pagestyle{fancy} % Finally, use the "fancy" page style to implement the FancyHdr headers

\newcommand{\HRule}{\rule{\linewidth}{0.5mm}} % New command to make the lines in the title page

\setlength{\parindent}{1cm}

% PDF meta-data
\hypersetup{pdftitle={\ttitle}}
\hypersetup{pdfsubject=\subjectname}
\hypersetup{pdfauthor=\authornames}
\hypersetup{pdfkeywords=\keywordnames}


%----------------------------------------------------------------------------------------
%   CAPA (ESTA PÁGINA É OBRIGATÓRIA)
%----------------------------------------------------------------------------------------


\begin{titlepage}

\noindent
\includegraphics[width=0.14\textwidth]{Imagens/logo_minerva_colorida.png}
\begin{minipage}[b]{0.73\textwidth}
\centering

{\fontsize{11pt}{12pt}\selectfont\textbf{
\univname\\
\campname\\
\deptname\\
% \facname\\ %% Descomentar para trabalhos do programa de pós-graduação
[4.5cm]
}}

\end{minipage}%
\includegraphics[width=0.16\textwidth]{Imagens/logo_valongo.png}

\vspace{4cm}

\begin{center}

{\fontsize{16pt}{18pt}\selectfont
\ttitle\\[6cm]
}
 
{\fontsize{12pt}{14pt}\selectfont
\authornames\\[3cm]
}

\end{center}

\begin{center}

\vfill

{\fontsize{12pt}{14pt}\selectfont 
\addressnames\\
\thesisdate %%%\today   %% comente esta linha e escreva a data no formato "Mês_por_extenso de ano" (ex: Agosto de 2022) caso não queira a data automática
}

\end{center}

\end{titlepage}

\clearpage % Start a new page


%----------------------------------------------------------------------------------------
%   FOLHA DE ROSTO (ESTA PÁGINA É OBRIGATÓRIA)
%----------------------------------------------------------------------------------------


\pagenumbering{arabic} % Begin numeric (1,2,3...) page numbering
\setcounter{page}{1}

\pagestyle{empty}

\begin{minipage}[b]{0.73\textwidth}
\centering
\hfill \hfill \hfill
\end{minipage}

\vspace{2cm}

\begin{center}

{\fontsize{16pt}{18pt}\selectfont
\ttitle\\[4cm]
}
 
{\fontsize{12pt}{14pt}\selectfont
\authornames\\[3cm]
}

\end{center}

\begin{textblock}{7}(9,10)

\noindent
{\fontsize{12pt}{14pt}\selectfont\justify{
\ttype \space submetido ao \deptname, \univname, como requisito necessário para a obtenção do título de \degreename. %% versão para trabalhos de graduação
% \ttype \space submetida ao \facname \space do \deptname, \univname, como requisito necessário para a obtenção do título de \degreename.   %% comente esta linha e descomente a linha abaixo se o seu trabalho estiver em inglês.
% \ttype submitted to the \facname of the \deptname, \univname, in fulfillment of the requirements for the degree of \degreename
%% ATENÇÃO: modificar para "como parte dos requisitos necessários para se o trabalho for de qualificação e em português.
%% ATENÇÃO: modificar para "in partial fulfillment of" se o trabalho for de qualificação e em inglês.
\\[1cm]
}}  

\noindent
{\fontsize{12pt}{14pt}\selectfont\justify{
Orientador: \supname\\[3cm]
}}

\end{textblock}

\begin{center}

\vfill

{\fontsize{12pt}{14pt}\selectfont 
\addressnames\\
\thesisdate %%%\today   %% comente esta linha e escreva a data no formato "Mês_por_extenso de ano" (ex: Agosto de 2022) caso não queira a data automática
}

\end{center}

\clearpage


%----------------------------------------------------------------------------------------
%   FICHA DE APROVAÇÃO (ESTA PÁGINA É OBRIGATÓRIA APÓS A DEFESA)
%----------------------------------------------------------------------------------------


\newpage
\thispagestyle{empty}
\mbox{}
\newpage


%----------------------------------------------------------------------------------------
%   DEDICATÓRIA (ESTA PÁGINA É OPCIONAL)
%----------------------------------------------------------------------------------------


\pagestyle{empty}

\setstretch{1.3} % Return the line spacing back to 1.3

% Page style needs to be empty for this page

\begin{center}

\input{Pre-Texto/Dedicatoria}

\end{center}

%\addtocontents{toc}{\vspace{2em}} % Add a gap in the Contents, for aesthetics

\clearpage


%----------------------------------------------------------------------------------------
%   AGRADECIMENTOS (ESTA PÁGINA É OPCIONAL)
%----------------------------------------------------------------------------------------


\setstretch{1.3} % Reset the line-spacing to 1.3 for body text (if it has changed)

\input{Pre-Texto/Agradecimentos}

\clearpage


%----------------------------------------------------------------------------------------
%   EPÍGRAFE (ESTA PÁGINA É OPCIONAL)
%----------------------------------------------------------------------------------------


\pagestyle{empty} % No headers or footers for the following pages

\null\vfill\vfill\vfill\vfill\vfill\vfill\vfill % Add some space to move the quote down the page a bit

\begin{flushright}
{\normalsize 

\input{Pre-Texto/Epigrafe}

}
\end{flushright}

\vfill\vfill\vfill\null % Add some space at the bottom to position the quote just right

\clearpage


%----------------------------------------------------------------------------------------
%   RESUMO (ESTA PÁGINA É OBRIGATÓRIA)
%----------------------------------------------------------------------------------------


%\addtotoc{Resumo} % Add the "Abstract" page entry to the Contents

%\addtocontents{toc}{\vspace{1em}} % Add a gap in the Contents, for aesthetics
\begin{resumo}

\input{Pre-Texto/Resumo}

\end{resumo}

\clearpage


%----------------------------------------------------------------------------------------
%   ABSTRACT (ESTA PÁGINA É OBRIGATÓRIA)
%----------------------------------------------------------------------------------------


%\addtotoc{Abstract} % Add the "Abstract" page entry to the Contents

%\addtocontents{toc}{\vspace{1em}} % Add a gap in the Contents, for aesthetics

\begin{abstract}

\input{Pre-Texto/Abstract}

\end{abstract}

\clearpage


%----------------------------------------------------------------------------------------
%   LISTA DE FIGURAS/LISTA DE TABELAS (ESTA PÁGINA É OBRIGATÓRIA)
%----------------------------------------------------------------------------------------


\pagestyle{fancy} % The page style headers have been "empty" all this time, now use the "fancy" headers as defined before to bring them back

\pagestyle{empty}
\lhead{\textit{Lista de Figuras}} % Set the left side page header to "List of Figures"
\listoffigures % Write out the List of Figures
\pagestyle{plain}

\lhead{\textit{Lista de Tabelas}} % Set the left side page header to "List of Tables"
\listoftables % Write out the List of Tables


%----------------------------------------------------------------------------------------
%   ABREVIAÇÕES (ESTA PÁGINA É OPCIONAL, MAS SUGIRO INCLUÍ-LA)
%----------------------------------------------------------------------------------------


\setstretch{1.5} % Set the line spacing to 1.5, this makes the following tables easier to read

\lhead{\textit{Abreviações}} % Set the left side page header to "Abbreviations"

\input{Pre-Texto/Abreviacoes}
\clearpage


%----------------------------------------------------------------------------------------
%   CONSTANTES FÍSICAS E OUTRAS DEFINIÇÕES (ESTA PÁGINA É OPCIONAL, MAS SUGIRO INCLUÍ-LA)
%----------------------------------------------------------------------------------------


\lhead{\textit{Constantes Físicas}} % Set the left side page header to "Physical Constants"

\input{Pre-Texto/Constantes}

\clearpage


%----------------------------------------------------------------------------------------
%   SÍMBOLOS (ESTA PÁGINA É OPCIONAL, MAS SUGIRO INCLUÍ-LA)
%----------------------------------------------------------------------------------------


\lhead{\textit{Símbolos}} % Set the left side page header to "Symbols"

\input{Pre-Texto/Simbolos}

\clearpage


% -----------------------------------------------------------------------------
%   SUMÁRIO (ESTA PÁGINA É OBRIGATÓRIA)
% ------------------------------------------------------------------------------------

\pagestyle{empty}
\lhead{\textit{Sumário}} % Set the left side page header to "Contents"
\tableofcontents % Write out the Table of Contents
\pagestyle{plain}

%----------------------------------------------------------------------------------------
%   CAPÍTULOS DA TESE
%----------------------------------------------------------------------------------------


\pagestyle{fancy} % Return the page headers back to the "fancy" style

\makeatletter   %% primeira página de cada capítulo contada, mas não numerada
\let\ps@plain\ps@empty   %% primeira página de cada capítulo contada, mas não numerada
\makeatother   %% primeira página de cada capítulo contada, mas não numerada

% Include the Capitulos of the thesis as separate files from the Capitulos folder

\input{Capitulos/Chapter1.tex}
\input{Capitulos/Chapter2.tex}
%%%você pode adicionar mais capítulos, se necessário, escrevendo uma nova linha no mesmo formato das linhas acima (ex: \input{Capitulos/Chapter6.tex}) e criando um novo arquivo de mesmo nome (ex: Chapter6) na pasta Capitulos

\clearpage


%----------------------------------------------------------------------------------------
%   BIBLIOGRAFIA (ESTA PÁGINA É OBRIGATÓRIA)
%----------------------------------------------------------------------------------------


\label{Bibliography}

\addtocontents{toc}{\vspace{1em}} % Add a gap in the Contents, for aesthetics

\addtotoc{Referências Bibliográficas}

\lhead{\emph{Bibliografia}} % Change the page header to say "Bibliography"

\bibliographystyle{Bibliografia/aasjournal_modified}  % Use the "unsrtnat" BibTeX style for formatting the Bibliography

\bibliography{Bibliografia/Bibliografia} % The references (bibliography) information are stored in the file named "Bibliography.bib"


%----------------------------------------------------------------------------------------
%   APÊNDICES (ESTAS PÁGINAS SÃO OPCIONAIS)
%----------------------------------------------------------------------------------------


\addtocontents{toc}{\vspace{1em}} % Add a gap in the Contents, for aesthetics

\appendix % Cue to tell LaTeX that the following 'Capitulos' are Apendices

% Include the appendices of the thesis as separate files from the Appendices folder

\input{Apendices/AppendixA}
%%% você pode adicionar mais apêndices, se necessário, escrevendo uma nova linha no mesmo formato das linhas acima (ex: \input{Apendices/AppendixC.tex}) e criando um novo arquivo de mesmo nome (ex: AppendixC) na pasta Apendices

%\addtocontents{toc}{\vspace{2em}} % Add a gap in the Contents, for aesthetics



\end{document}  